\documentclass[11pt]{article}
\usepackage{fontspec}
\setmainfont{Times New Roman}
\setsansfont{Arial}
\setmonofont{Consolas}
\usepackage{amsmath,amssymb,mathtools}
\usepackage{geometry}
\usepackage{microtype}
\geometry{a4paper,margin=2.35cm}
\setlength{\parindent}{1.25em}
\setlength{\parskip}{0pt}
\emergencystretch=25em
\sloppy
\hbadness=10000
\vbadness=10000
\hfuzz=2pt
\vfuzz=10pt
\raggedbottom

\begin{document}

% Unit: Paper27 v001. Direct Interslavic Latin translation from the German final-audited source slice.
\section*{27. Hilbertove čisla v teoriji idealov}

\begin{center}
J. Ber. d. DMV 34 (1925), str. 101
\end{center}

E. Noether: Hilbertove čisla v teoriji idealov. Pod Hilbertovymi čislami obče razumějut se taka čisla, ktore odnoset se k izčisljenju ostatkovyh klasov idealov. Sam Hilbert v svojoj „karakterističnoj funkciji“ dal funkciju, ktora za idealy polinomnoj oblasti od někojej granicy dalje dava točno čislo linearno nezavisnyh ostatkovyh klasov. Za nižše stepenove čisla Macaulay dopolnil tu funkciju tvorjačoju funkcijeju, ktoru Ostrowski mogl eksplicitno izračunati. Ale Macaulay je takože podal čisla drugogo vida, ktore se pokazali kako najbitnějše za obču teoriju idealov, ibo možno jih shvatiti abstraktno. Soglasno obščim razložnym teoremam dostatočno jest dati definiciju za primarne idealy $\mathfrak q$. Ide o čislo linearno nezavisnyh ostatkovyh klasov vzhodno k polju ostatkovyh klasov asociovanogo prostogo ideala $\mathfrak p$. Ta čisla razpadajut se na podsistemy, ktore vsakokrat sut dane čislom nerazložimyh komponentov idealov $\mathfrak q,\mathfrak q/\mathfrak p,\mathfrak q/\mathfrak p^2,\ldots$. Čislo $i$-togo podsistema jest takože karakterizovano čislom členov kompozicijnogo reda, ktory ide od $\mathfrak q/\mathfrak p^i$ k $\mathfrak q/\mathfrak p^{i-1}$. Cěly kompozicijny red i jego Hilbertove čisla tvorjat ekvivalent predstavjenja primarnyh idealov kako potencij prostyh idealov, ktoro v obščem slučaju -- napr. v konečnyh porjadkah algebraičnogo čislovogo polja -- veče ne važi.

\end{document}
